% No preamble. Prefix vst-. Actual composite-constraint vector UV bound.
\section{A uniform short-time estimate for the actual vector operator}
\label{vst-section}

Fix the primitive odd ratio \(\ell\ge3\), the actual selected analytic
source mean, and \(a>0\). Let \(K=\ell^k\), and let \(M\) be any
preceding homogeneous depth. We work in an actual joint infinite
structural system, with current-link angle counting coordinates.
All affine substitutions below are polynomial jets at zero.
No continuously varying periodic magnetic flux is used.

Choose \(\mu\ne\nu\) and a Cartan element \(H\) with
\(\tau(H^2)=1\), and write
\[
 U_\mu(t;x)=\exp(i t x_\nu H/K^2),\qquad U_\lambda(t;x)=1
                       \quad(\lambda\ne\mu).
\]
The stationary-jet identification established in Notebook 33 gives,
through order two, the actual vector family
\begin{equation}
 \begin{split}
 B_M(t)&=s_M''(U(t))+D_b(t)D_b(t)^T,\\
 V_{M,K}(t)&=B_M(t)+aC_K(t)^TC_K(t).
 \end{split}
 \label{vst-operators}
\end{equation}
Here \(s_M''\) is the actual left-relative Hessian,
\(D_b\) has scalar values zero at the \(K\)-centers, and \(C_K\)
is the left-relative derivative of the composite ordinary map
\(M_\ell^k\). In particular, \(s_M''\) is not replaced by the
bare Wilson Hessian when \(M>1\).

The inputs used from the classical construction are uniform
current-coordinate exponential bounds for \(q_M=s_M''(1)\),
the complete scalar-action support/Cauchy bounds through order
four on a fixed complex product ball, and the exact local gauge
covariance of the relative Hessian. The local coordinate and
color fibers have fixed dimension. The actual primitive
averaging formula supplies finite-range analytic derivative
matrices. Its composite output row reads only the two final
endpoint blocks. These inputs are uniform in preceding depth
and compatible regulator volume.

\begin{theorem}[Actual vector ultraviolet bound]
\label{vst-theorem}
There is a constant \(C\), depending on the fixed primitive
data, \(a\), the classical analytic bounds and the color group,
but not on \(M,K\) or the selected structural system, such that
\begin{equation}
 \int_0^1\left|[t^2]\operatorname{Tr}_{0,E}
                         e^{-sV_{M,K}(t)}\right|\frac{ds}{s}
 \le C.
 \label{vst-main}
\end{equation}
The trace is over fine bonds with initial site in block zero,
including color. Together with the actual ghost ultraviolet
bound, this proves the full short-time estimate on
\(0<s<1\), or \(0<u<K^{-2}\) after \(s=K^2u\).
\end{theorem}

\subsection{Uniform flat bounds for the full composite constraint}
\label{vst-flat-section}

Let \(D_r\) be the scalar forward difference on the lattice
of spacing \(r\), without division by \(r\). Define the
straight angle average
\begin{equation}
 (A_r v)(y,\lambda)=r^{-d}
       \sum_{x\in B_r(y)}\sum_{h=0}^{r-1}
                              v_\lambda(x+h e_\lambda).
 \label{vst-straight}
\end{equation}
Let \(b_\ell\) be the scalar mean of internal primitive
root-contour sums. The actual arithmetic differential is
\begin{equation}
 C_\ell=A_\ell-D_\ell b_\ell,\qquad
 C_\ell D_1=D_\ell R_\ell,
 \label{vst-primitive}
\end{equation}
where \(R_\ell\) restricts scalar values to the new centers.
Tiling centered blocks and concatenating strings gives
\(A_\ell^{(r)}A_r=A_{\ell r}\). Consequently
\begin{equation}
 \begin{split}
 C_K&=A_K-D_Kb_K,\\
 b_{\ell r}&=b_\ell^{(r)}A_r+R_\ell b_r,\\
 b_K&=\sum_{j=0}^{k-1}
       R_{\ell^{j+1}\to K}\,b_\ell^{(\ell^j)}A_{\ell^j}.
 \end{split}
 \label{vst-recursion}
\end{equation}
Each scalar restriction in the last line is a counting
\(L^2\) contraction. At a final center it merely retains
the same nested center; it does not sum all intermediate
centers.

The absolute row sums of \(A_r\) equal \(r\), and its
absolute column sums equal \(r^{1-d}\). For a column,
the two neighboring longitudinal output intervals have
complementary multiplicities summing to \(r\), and the
transverse output block is unique. Thus
\begin{equation}
                \|A_r\|_{2\to2}\le r^{1-d/2}.
 \label{vst-straight-norm}
\end{equation}
These counting identities retain their multiplicities
on compatible periodic graphs. Since \(b_\ell\) is a
fixed finite-range operator and
\(\|D_K\|_{2\to2}\le2\sqrt d\), in \(d=4\)
\begin{equation}
 \|C_K\|_{2\to2}
 \le K^{-1}+
       \frac{2\sqrt d\,\|b_\ell\|_{2\to2}}{1-\ell^{-1}}.
 \label{vst-full-flat}
\end{equation}
Every consecutive flat composite segment therefore has
norm at most one fixed \(C_0\ge1\), including empty segments.
This is an \(O(1)\) bound for the full ordinary map, not
the generally false \(O(K^{-1})\) assertion.

For localized information, a row of \(A_r\) has at most
\(2r^d\) entries, each at most \(r^{1-d}\), so
\[
              \|A_r(c,\cdot)\|_2\le\sqrt2\,r^{-1}
                         \quad(d=4).
 \]
Equation \eqref{vst-recursion} gives uniformly bounded
rows of \(b_K\), and hence of \(C_K\).
The two-block dependency also shows that each input
bond occurs in only a dimension-dependent number of
different final output rows. This gives a second direct
row-and-incidence proof of the full bound.

\subsection{Actual affine jets of the nonlinear composite derivative}
\label{vst-magnetic-section}

Fix first a final row with endpoint blocks near zero.
At stage \(j\), put \(r=\ell^j\). The exact affine
averaging identity gives the intermediate background
angle, in that stage's lattice coordinates,
\[
                     w_j(z)=t(r^2/K^2)z_\nu H.
\]
The exact nested dependency geometry confines every
needed coordinate to \(|z|\le c_\ell K/r\). Therefore
\begin{equation}
                  \|w_j\|_{\rm needed}
                            \le c_\ell |t|r/K.
 \label{vst-stage-smallness}
\end{equation}
At each stage independently, extend the needed values
to a globally bounded field with this same bound, by
clipping the coordinate outside a larger dependency box.
The extensions need not be successive blockings.
Every derivative entry contributing to the selected
product row agrees with the actual affine stage, which
is the only use of these extensions.

Let \(F_j\) be the flat primitive derivative and
\(E_j(t)\) its change at the extended background.
The fixed primitive analytic and incidence bounds
give, on one disk \(|t|<\rho_0\),
\begin{equation}
                  \|E_j(t)\|_{2\to2}
                          \le C_\ell |t|\ell^j/K.
 \label{vst-stage-error}
\end{equation}
Expand the derivative product by the subset of stages
where \(E_j\) is used. A term with \(n\) marked stages
has \(n+1\) intervening flat segments, each bounded by
\(C_0\). Hence
\begin{align}
 \left\|\prod_j(F_j+E_j(t))\right\|
 &\le C_0\prod_j(1+C_0\|E_j(t)\|)\notag\\
 &\le C_0\exp\!\left(C C_0|t|
                         \sum_{j<k}\ell^j/K\right)\notag\\
 &\le C_0\exp(C C_0|t|/(\ell-1)).
 \label{vst-product-bound}
\end{align}
The actual selected row consequently has a uniform
analytic \(L^2\) bound. Cauchy estimates on a smaller
fixed disk bound its Taylor coefficients through order two.

For a row at coarse position \(y\), translating the
affine field adds a constant Cartan potential. Remove
it by the site gauge \(\exp(i t y_\nu x_\mu H/K)\).
Exact gauge covariance expresses the row using its
centered counterpart and input/output adjoint frames.
On its two \(K\)-blocks the frame derivatives through
order two are bounded by a fixed polynomial in \(|y|\),
independently of \(K\). Writing
\(C_K(t)=\sum_j t^j C^{[j]}\), one may take
\begin{equation}
        \|C^{[j]}(y,\cdot)\|_2
                  \le C_j(1+|y|)^{2j},
                  \qquad j=0,1,2.
 \label{vst-row-jets}
\end{equation}
Their support remains in the same two endpoint blocks.
This is a statement about polynomial jets; it does not
assert a bounded analytic operator at a nonzero unbounded
affine field.

Group current bonds by their initial-site \(K\)-block,
and final bonds by their initial coarse site. The latter
fiber has fixed dimension \(r_c=d\dim\mathfrak g\).
Equation \eqref{vst-row-jets} and finite incidence give
finite-range block operator bounds with polynomial
growth for \(C^{[j]}\) and their transposes.
In particular \(C^{[0]}\) has a uniform weighted block
Schur norm for every fixed block exponent.

\subsection{The base vector before the constraint insertion}
\label{vst-base-section}

At zero,
\[
                 B_0=q_M+D_bD_b^T
\]
is positive and has a uniform exponentially weighted
fine-lattice Schur norm. A spectral gap is unnecessary
on \(0\le s\le1\): the Schur-algebra power series bounds
its semigroup uniformly with exponential matrix decay.

The actual scalar-action bounds through order four
give uniform exponentially rooted kernels for the
first two background derivatives of \(B\).
The nearest-neighbor based-gradient term has the same
property uniformly in the center spacing.
For each fine diagonal bond rooted at \(x\), use a
Cartan gauge to subtract the constant potential
\(t x_\nu H/K^2\). The diagonal color heat trace is
unchanged. Its source field is now
\(t(z_\nu-x_\nu)H/K^2\). Each linear magnetic vertex
has a factor \(K^{-2}\), and each quadratic vertex
has \(K^{-4}\); uniform exponential kernels absorb
the remaining coordinate polynomials.

The one quadratic insertion and the ordered two linear
insertions in the heat expansion therefore give
\begin{equation}
 \left|[t^2]\operatorname{tr}_{\mathfrak g}
                         (e^{-sB(t)})_{xx}\right|
        \le CK^{-4}(s+s^2)e^{Cs},\qquad 0<s\le1.
 \label{vst-base-diagonal}
\end{equation}
The local gauge identities extend to these polynomial
jets by their exponentially summable kernels.
There are \(dK^4\) bond roots per block, so
\begin{equation}
 \left|[t^2]\operatorname{Tr}_{0,E}e^{-sB(t)}\right|
                  \le C(s+s^2)e^{Cs}.
 \label{vst-base-trace}
\end{equation}
Its \(ds/s\) integral is uniformly finite.

\subsection{Uniform block estimates without tracing the large fiber}
\label{vst-block-section}

The coefficients \(B^{[j]}\), \(j\le2\), also have
uniform exponentially decaying \(K\)-block operator
kernels, with fixed polynomial growth in the block
indices for \(j>0\). Fine positions in block \(X\)
have magnitude at most \(CK(1+|X|)\), while the
magnetic source coefficients contain \(K^{-2j}\).
Fine Schur bounds control the operator norm of each
block without its dimension as a prefactor.
For distant blocks the fine distance is at least
\(K(d_{\rm block}(X,Y)-d)\); a smaller fixed block
exponent is therefore summable uniformly in \(K\).
The finitely many nearby block offsets are controlled
by the full fine operator norm.

Set \(A(t)=aC_K(t)^TC_K(t)\). The preceding finite-range
block bounds apply to \(A^{[j]}\) and to
\(V^{[j]}=B^{[j]}+A^{[j]}\).
For \(E=B,V\), the block power series at zero and the
one/two insertion formulas imply, for \(j\le2\),
\begin{equation}
 \big\|([t^j]e^{-sE(t)})_{XY}\big\|
 \le C(1+|X|+|Y|)^p e^{-\kappa|X-Y|},
                         \qquad 0\le s\le1.
 \label{vst-semigroup-block}
\end{equation}
All constants are uniform in \(M,K\).
For positive \(j\) the usual factors \(s\) and
\(s+s^2\) can also be retained. Products of polynomial
kernels are controlled by spending a smaller block
exponent. No trace over a \(K^4\)-dimensional fiber
has yet been taken.

\subsection{The constraint correction has fixed coarse trace rank}
\label{vst-rank-section}

Let \(P_0\) project onto current bonds rooted in block zero.
The exact Duhamel difference is
\begin{equation}
 e^{-sV(t)}-e^{-sB(t)}
   =-a\int_0^s e^{-(s-u)V(t)}
              C_K(t)^TC_K(t)e^{-uB(t)}\,du.
 \label{vst-difference}
\end{equation}
At each coefficient through degree two, cyclicity gives
\begin{equation}
 \operatorname{Tr}_{\rm fine}
           (P_0E_VC^TCE_B)
       =\operatorname{Tr}_{\rm coarse}
           (CE_BP_0E_VC^T).
 \label{vst-local-cyclicity}
\end{equation}
The projector \(P_0\) remains in this identity. The
right side sums all coarse diagonal blocks; it is
not identified with an arbitrary block-zero coarse
trace density.

For each coarse root \(Y\), its diagonal block on the
right is a product through the fine block-zero Hilbert
space. The outer fiber has dimension \(r_c\), independent
of \(K\). Equations \eqref{vst-row-jets} and
\eqref{vst-semigroup-block} bound the coefficients through
degree two of the two maps by
\[
                     C(1+|Y|)^p e^{-\kappa'|Y|}.
\]
The absolute trace of their product is at most \(r_c\)
times the product of these norms. Summing \(Y\) is
uniformly finite. Expanding the quadratic coefficient
in \eqref{vst-local-cyclicity}, and then integrating
the Duhamel time, proves
\begin{equation}
 \left|[t^2]\operatorname{Tr}_{0,E}
          (e^{-sV(t)}-e^{-sB(t)})\right|
                    \le Cs,\qquad 0<s\le1.
 \label{vst-rank-bound}
\end{equation}
One can first use finite boxes and compact sources.
The displayed rooted sums are absolute and justify
the limit. For fixed \(K\), the finite \(P_0\) also
makes ordinary trace cyclicity legitimate; the
estimate itself uses the small outer coarse fiber
to avoid a \(K^4\) factor. Complex coefficients use
the holomorphic bilinear transpose, which agrees
with the adjoint at the real flat point.

Combining \eqref{vst-rank-bound} with
\eqref{vst-base-trace} and integrating \(ds/s\)
proves Theorem~\ref{vst-theorem}.

\subsection{What this estimate does and does not settle}
\label{vst-scope}

The argument uses the actual composite fixed-\(\ell\)
hierarchy. A direct-\(K\) contour average can have
different raw row bounds and was not substituted.
The full ordinary map needs only the proved
\(O(1)\) counting bound; sampled-center gauge
directions forbid assuming \(O(K^{-1})\) for it.

The final affine statement is in the infinite
structural system. The uniform finite-regulator
inputs and local limits do not turn it into a
continuous magnetic-flux formula on an arbitrary
torus. The semigroup estimates here use \(e^{Cs}\)
on bounded time, not a uniform large-time gap.
They prove neither the continuum coefficient on
intermediate times nor the cancellation of the
soft based-ghost contribution at large times.
A combined signed-integral comparison remains
an independent problem.
